% !TEX root = ../main.tex
\chapter{O editor}
\label{cap:editor}

O editor da AURORA é construído sobre o \tool{Monaco}, o mesmo motor de edição do Visual Studio Code. Multi-cursor, minimapa, localizar e substituir, desfazer ilimitado e a ergonomia geral que você já conhece vêm de fábrica; este capítulo cobre o que a AURORA acrescenta por cima.

\section{Abas}

Cada arquivo aberto vira uma aba na barra superior do editor, com um ponto indicando alterações não salvas. As abas podem ser reordenadas por arraste. Salva-se a aba atual com \tecla{Ctrl+S} e todas com \tecla{Ctrl+Shift+S}; fecha-se a atual com \tecla{Ctrl+W}, reabre-se a última fechada com \tecla{Ctrl+Shift+T}, e \tecla{Ctrl+N} cria um arquivo novo.

\section{Editor dividido}
\label{sec:split}

O botão de divisão, que flutua no canto do painel em foco, abre o arquivo atual num novo painel lado a lado, até o limite de três painéis, cada um com a sua barra de abas. O detalhe importante é que painéis exibindo o mesmo arquivo compartilham o mesmo documento: editar de um lado atualiza o outro ao vivo, com desfazer e estado de salvamento unificados. É a forma natural de manter o \file{.cmm} à esquerda e o \file{.v} gerado à direita, ou duas regiões distantes do mesmo arquivo.

\figurapendente{Editor dividido em dois painéis}{O media\_movel.cmm à esquerda e o media\_movel.v gerado à direita.}

\section{Realce de sintaxe}

A \cmm{} tem gramática de realce própria, cobrindo palavras-chave, diretivas, funções da biblioteca, literais complexos e a notação de Dirac; os \code{\#define} do programa são recolhidos dinamicamente e coloridos como constantes. O \textit{assembly} do SAPHO também tem gramática própria. O Verilog e o SystemVerilog usam a gramática do Monaco, enriquecida com \textit{tokens} semânticos do tree-sitter, que igualmente atende C e C++ no caminho alternativo de compilação.

\section{Diagnósticos e navegação}

Dois servidores de linguagem (LSP, \textit{Language Server Protocol}) trabalham nos arquivos Verilog. O \tool{Verible} é o servidor completo: diagnósticos de sintaxe e estilo enquanto se digita, formatação, \textit{outline} de símbolos, \textit{hover}, ir à definição e referências. O \tool{slang} complementa com a análise semântica do projeto inteiro, elaborando todos os módulos em conjunto para apontar erros de elaboração e portas incompatíveis, além de oferecer autocompletar; por ser mais pesado, pode ser ligado e desligado com \tecla{Ctrl+Alt+S}. Os problemas aparecem sublinhados no código, e tudo funciona em regime de melhor esforço: sem um servidor disponível, a edição segue sem ele.

\section{Formatação automática}

O atalho \tecla{Shift+Alt+F} formata o documento atual, para C, C++ e \cmm{}, com o \tool{clang-format} empacotado. Se o projeto tiver um arquivo \file{.clang-format} próprio, ele é respeitado.

\section{Busca no projeto}

O atalho \tecla{Ctrl+Shift+F}, ou o botão de busca no cabeçalho da árvore, abre o modal Buscar nos Arquivos: termo com opções de diferenciar maiúsculas, palavra inteira e expressão regular, resultados agrupados por arquivo e clique levando à linha. A busca ignora o \file{.git} e as pastas de artefatos.

\section{Seleção e IA}

Ao selecionar um trecho de código, um pequeno botão em forma de estrela aparece junto à seleção: ele envia o trecho diretamente à Aurora Intelligence, já contextualizado. É o caminho mais curto para pedir uma explicação ou perguntar por que algo não compila (Capítulo~\ref{cap:aurora-intelligence}).

\section{Conflitos com o disco}

Se um arquivo aberto for alterado por outro programa, a AURORA detecta a divergência e pergunta o que fazer no diálogo Arquivo Modificado Externamente, oferecendo manter a versão do editor, usar a versão do disco ou salvar e recarregar. Nada é sobrescrito sem a sua decisão.

\section{Atalhos de teclado}

Os atalhos citados neste capítulo, todos reconfiguráveis nas Configurações, estão reunidos com os demais no Apêndice~\ref{ap:atalhos}. Dentro do editor valem ainda os atalhos usuais do Monaco, como \tecla{Ctrl+F} para localizar, \tecla{Alt+Click} para múltiplos cursores e \tecla{Ctrl+/} para comentar a linha. A posição do cursor é exibida na barra de status.
